% Chip Scanner: Label-Free Optical Defect Inspection and Metrology for Patterned Wafers
% Build:  cd paper && latexmk -pdf chipscan.tex
% Figures: uv run python scripts/paper/make_figures.py   (writes paper/figures/)
\documentclass[journal]{IEEEtran}

\usepackage[T1]{fontenc}
\usepackage{textcomp}
\usepackage{amsmath,amssymb}
\usepackage{graphicx}
\usepackage{booktabs}
\usepackage{array}
\usepackage{url}
\usepackage[hidelinks]{hyperref}
\usepackage{cite}

\graphicspath{{figures/}}
\newcommand{\um}{\,\textmu m}
\newcommand{\umpx}{\,\textmu m/px}

\begin{document}

\title{Chip Scanner: Label-Free Optical Defect Inspection and Metrology for Patterned Wafers}

\author{Clay~Nicholson%
\thanks{C. Nicholson is with the Research Laboratory of Electronics, Massachusetts Institute of Technology, Cambridge, MA 02139 USA.}}

\markboth{}{Nicholson: Label-Free Optical Defect Inspection and Metrology for Patterned Wafers}

\maketitle

\begin{abstract}
Wafer-scale optical inspection outside high-volume fabrication is still largely manual: a reviewer pans a gigapixel brightfield scan and judges each particle by eye, because the die-to-die and die-to-database comparison of industrial inspection tools assumes many identical dies, controlled imaging and labelled defect libraries. We present Chip Scanner, a pipeline that converts a stitched optical scan of any patterned chip or wafer into a ranked list of findings, each assigned a likely failure mode and a risk score from 0 to 100. Preprocessing corrects per-frame illumination and tilt and restricts processing to the patterned die. Four detectors are fused: a speck segmentation network, a design-free anomaly network, a twin-tile reference residual and a pad rule. A verification chain driven by a layout segmentation network then removes candidates that belong to the design. All networks are trained without manual labels, from defects injected into clean tiles of the target scans, so the system adapts to a new process from its own images. Only the risk model encodes a device technology. We evaluate the system on superconducting quantum circuit wafers, whose micrometre-scale defects and large repeated arrays make a demanding test vehicle. On a held-out chip with 300 planted specks, the full pipeline recovers 93.3\% of them (86.7\% of those below 12~px) at 1.2 detections per clean $768 \times 768$~px window. A 100-mm wafer of $118{,}092 \times 118{,}088$~px is processed in 37~min on a laptop and yields 119{,}122 findings, of which 2{,}064 are rated medium risk or higher. When the design file is available, the scan is registered to it with a median residual of 4.2~px, and line and gap widths are measured at 16{,}222 sites. They are reproducible to 0.1\um{} across two wafers of the same design.
\end{abstract}

\begin{IEEEkeywords}
Automated optical inspection, wafer inspection, defect detection, semantic segmentation, synthetic training data, critical-dimension metrology, superconducting quantum circuits.
\end{IEEEkeywords}

\section{Introduction}
\IEEEPARstart{P}{articles}, residues, scratches and pattern errors at the micrometre scale decide the yield of almost every patterned device, from CMOS and MEMS to integrated photonics and superconducting circuits. Brightfield optical microscopy is the fastest way to screen for them, but its output is large. At 1.6\umpx, one 100-mm wafer is a $118{,}092 \times 118{,}088$~px image, about $1.4 \times 10^{10}$ pixels, and a reviewer panning it by eye covers only a small fraction.

Industrial automated optical inspection (AOI) solves this with die-to-die or die-to-database comparison under tightly controlled imaging, backed by defect libraries built up over many production lots. Research, prototyping and low-volume fabrication meet none of these conditions. Wafers carry several reticle designs in small numbers, scans are stitched from thousands of camera frames with per-frame illumination falloff and sub-pixel jumps, and no labelled defect data exist for a new process. Unsupervised anomaly detection~\cite{bergmann2019mvtec} and training on synthetic anomalies~\cite{li2021cutpaste,zavrtanik2021draem} address the label problem in general imagery, but not the registration, scale and risk-ranking problems of a wafer scan.

This paper describes Chip Scanner, a process-agnostic system for this setting. Its detection, verification, registration and metrology stages use no knowledge of the device technology. They learn from the target scans and, where available, from the GDS layout. Only a small risk table maps a finding's location to the failure modes of a given technology. We use superconducting quantum circuit wafers from the MIT Research Laboratory of Electronics as the test vehicle. Their coplanar-waveguide (CPW) resonators and Josephson junctions fail from defects of a few micrometres: a particle bridging a CPW gap shorts the line, residue in a resonator gap adds two-level-system dielectric loss, and debris on a junction alters or destroys it~\cite{krantz2019guide,muller2019tls}. Their large repeated arrays and stitched Olympus DSX1000 scans (1.0--1.6\umpx) exercise every stage of the pipeline. The contributions are as follows:
\begin{enumerate}
\item A process-agnostic pipeline from a stitched scan to ranked, failure-mode-labelled findings, running on a laptop at wafer scale.
\item Label-free training of all detection networks by defect injection into clean tiles of the target scans, supplemented by real contamination harvested from twin-tile disagreement, so that a new process needs no annotation.
\item A candidate verification chain whose layout-dependent rejections are conditioned on local context; on one wafer, conditioning recovered 94{,}600 real particles that a context-free rule discarded.
\item Wafer-scale registration to the GDS layout and sub-pixel line and gap metrology at more than 16{,}000 sites per wafer.
\item An evaluation protocol on a held-out chip with planted defects, and a reported negative result for device-to-device comparison.
\end{enumerate}

Section~\ref{sec:overview} gives an overview of the system and Section~\ref{sec:methods} details its methods. Section~\ref{sec:training} covers training. Section~\ref{sec:results} reports results on the superconducting test vehicle, Section~\ref{sec:discussion} discusses generality and limitations, and Section~\ref{sec:conclusion} concludes.

\section{System Overview}\label{sec:overview}
Fig.~\ref{fig:pipeline} shows the pipeline; no stage before the risk model depends on the device technology. A scan enters as an Olympus \texttt{.vsi} file and is converted once into a multi-resolution zarr pyramid, from which every later stage and the review interface read. Preprocessing removes per-frame illumination falloff, corrects tilt and computes a chip mask. Four networks then run in a single fused pass over 1{,}024~px tiles: a layout network that labels traces, devices and lettering, and three detectors (speck, anomaly and trace). Together with a reference residual and a pad rule, the detectors propose candidate regions. A verification chain rejects candidates that are part of the design, and a risk model assigns each survivor a failure mode and a score. When a GDS layout is available, an optional branch registers the scan to it and measures line and gap widths. Other microscope formats enter through the same pyramid conversion.

\begin{figure*}[!t]
\centering
\includegraphics[width=\textwidth]{fig1_pipeline.png}
\caption{Processing pipeline. The upper path runs on every scan; the lower path runs only when the design file is available.}
\label{fig:pipeline}
\end{figure*}

All stages run in one process on a laptop (Apple M5) or a workstation GPU. Tiles are read ahead and written behind on background threads, and networks run in half precision, so inference is not limited by storage. Findings, intermediate maps and the reasons for every rejected candidate are stored alongside the scan, and a local web interface built on OpenSeadragon~\cite{openseadragon} streams the pyramid for review.

\section{Methods}\label{sec:methods}

\subsection{Preprocessing}
The DSX1000 stitches a scan from camera frames, $1{,}400 \times 1{,}043$~px on the NFT1 wafer. The frame grid is recovered from periodic brightness steps, a flat-field gain is estimated per frame position and divided out, and the scan is rotated by the angle that straightens the chip edges. A chip mask is computed on a level at most 4{,}096~px wide: patterned pixels are found by local contrast, gaps of flat ground plane between them are bridged by morphological closing, holes are filled and a 400\um{} margin is added. Stage, fill and bare wafer rim outside the mask are never processed or shown for review.

\subsection{Candidate Detection}
Four detectors produce per-pixel maps, and the union of their supra-threshold regions forms the candidate set, because their misses are largely disjoint.
\begin{enumerate}
\item \emph{Speck network.} A U-Net~\cite{ronneberger2015unet} with 48 base channels segments particles, residue and scratches down to a few pixels.
\item \emph{Anomaly network.} A second U-Net is trained without design input to flag any region that does not resemble normal chip surface, which catches fibres, stains and lifted metal absent from the speck training distribution.
\item \emph{Reference residual.} For each tile, twin tiles with normalized cross-correlation $\geq 0.8$ are located elsewhere on the chip and aligned. A per-pixel minimum--maximum envelope is built over the twins, and pixels outside it by more than three times the tile noise are flagged. Metal edges are excluded because stitching displaces them by up to one pixel.
\item \emph{Pad rule.} A classical contrast test on large bond and probe pads detects dirt and probe scrubs on flat metal.
\end{enumerate}
A layout network (Section~\ref{sec:training}) labels every pixel as background, trace, device or lettering in the same pass.

\subsection{Candidate Verification}\label{sec:verification}
Each connected candidate passes an ordered chain of checks; every rejection is stored with its reason. The chain tests, in order, detector attribution; lettering; the scan border; isolation of reference-only candidates; device membership; twin verification, which rejects a candidate also present at the corresponding position of its twins; and growth to full extent on pads, where probe-needle scrubs are labelled separately.

The lettering and device checks rely on layout labels, and the layout network also assigns those labels to isolated particles. Both checks therefore reject a candidate only if it is not isolated: there must be patterned structure within its neighbourhood, or the connected region that differs from background must extend beyond the candidate by more than $12 + 0.5e$~px, where $e$ is the candidate extent. Candidates on open ground are always retained, and twin verification still removes those that belong to the design. The checks run in parallel across CPU cores, which reduced this stage's runtime by a factor of 9.3 with identical output.

\subsection{Risk Model}
For each finding, the local geometry is measured from the image: whether the finding lies on metal or in a gap, and the gap width, from the dark edge lines on either side. Each failure mode $m$ in Table~\ref{tab:risk} receives a weight $w_m$, and the risk score is
\begin{equation}
R = 100 \cdot \max_{m} w_m \cdot (0.7 + 0.3\,c) \cdot k_{\mathrm{pad}},
\end{equation}
where $c \in [0, 1]$ is the detector confidence and $k_{\mathrm{pad}}$, the pad factor, is 0.6 for probe marks on pads and 1 otherwise. Scores are binned into critical, high, medium and low for review.

\begin{table}[!t]
\caption{Failure-Mode Weights for Superconducting Circuits}
\label{tab:risk}
\centering
\begin{tabular}{@{}l>{\raggedright\arraybackslash}p{0.33\columnwidth}>{\raggedright\arraybackslash}p{0.36\columnwidth}@{}}
\toprule
Mode & Condition & Weight $w_m$ \\
\midrule
Short & In a gap and spanning $\geq 80\%$ of its width & 0.95 \\
Junction & Near a Josephson junction & $0.35$--$1.0 \times$ proximity, rising with size \\
Loss & On or beside a trace & 0.55 in a resonator gap, otherwise 0.25; $\times$ size factor \\
Cosmetic & All other findings & $0.12 \times$ size factor \\
\bottomrule
\end{tabular}
\end{table}

The risk table is the only technology-specific component of the pipeline. Its inputs (whether a finding lies on a line or in a gap, how much of the gap it spans, its distance to device structures and its size) are measured identically for any layout. Table~\ref{tab:risk} is the instance for superconducting circuits used in this work. Another process substitutes its own modes and weights, such as bridging between adjacent interconnect lines or particles on optical waveguides.

\subsection{Design Registration and Line-Width Metrology}
Wafer layouts are too large to flatten: the NFT1 test wafer, for example, contains three reticle designs, each placed 12 times; flattening it would yield about $10^{7}$ polygons. Each reticle is instead loaded once with a bounding-box index and rasterized only where a window requires it. Registration proceeds in three stages: a global affine transform (which recovers 1.6285\umpx), a smooth displacement field for slow drift, and a translation per camera frame for stitching jumps. The per-frame stage reduces the median residual from 13.3 to 4.2~px.

Line and gap widths are measured on cuts across each patterned line every 500--2{,}000\um; on the test vehicle these are the CPWs. In brightfield, every metal edge images as a thin dark line. Its position is refined to sub-pixel precision by fitting a parabola to the intensity minimum $I_0$ and its neighbours $I_{-1}$, $I_{+1}$:
\begin{equation}
\delta = \frac{I_{-1} - I_{+1}}{2\,(I_{-1} - 2 I_0 + I_{+1})}.
\end{equation}
The designed edge sequence is then slid along the cut to the best match, and each line and gap width is the distance between its two matched edges. Features narrower than 3~px are not reported.

\section{Training}\label{sec:training}

\subsection{Supervision Without Manual Labels}
No labelled defect set exists for these chips, so every network is trained on labels generated from the scans themselves. Training tiles are drawn from four scenes (chips A2, B2 and an AHQ1 reticle, plus a batch of printed-circuit scans); chip C1 is held out entirely for validation and for the benchmark of Section~\ref{sec:results}.

For the speck and anomaly networks, synthetic particles, fibres and smudges are painted into clean tiles in the manner of synthetic-anomaly training~\cite{li2021cutpaste,zavrtanik2021draem}. Their brightness is sampled from the same scan, so injected and real dirt share a tone distribution, and the painter's own masks serve as exact labels. Real contamination is added by harvesting reference-residual detections, which are present in one copy of a repeated structure and absent from its twins; these supplied 4{,}554, 3{,}512 and 870 labelled regions on the three chip scenes. For the layout network, a classical segmentation gives pseudo-labels for background, trace and device, and lettering is painted in with its own label. Class weights of 0.3, 1.0, 3.0 and 2.0 offset the dominance of background.

\subsection{Optimization}
Binary detectors minimize a weighted binary cross-entropy (positive weight 3) plus a soft Dice term. Optimization uses AdamW (weight decay $10^{-4}$) with a one-cycle learning-rate schedule peaking at $2 \times 10^{-3}$. The production speck model, speck\_v4, is a U-Net with 48 base channels trained on 40{,}000 tiles for 30 epochs; this took 54~min on an NVIDIA RTX PRO 6000 workstation, about 15 times faster per tile than the laptop.

\subsection{Training Curves}
Fig.~\ref{fig:training} reports validation metrics after each epoch. These tiles follow the training recipe, so the curves measure how well each network learned its task rather than its accuracy on a real scan; Section~\ref{sec:results} measures the latter. Speck precision exceeds 0.90 by epoch 4, while recall needs longer and collapses once, to 0.45 at epoch 10, before recovering the following epoch. Over the last ten epochs recall holds at 0.96--0.97, and the share of clean tiles with any false detection falls from 0.54 after the first epoch to about 0.10. For the layout network, lettering is learned last (IoU 0.01 after one epoch and 0.69 after ten), consistent with its small size and chip-to-chip variability.

\begin{figure*}[!t]
\centering
\includegraphics[width=\textwidth]{fig2_training.png}
\caption{Validation metrics per epoch. (a) Speck network speck\_v4: recall, precision and the fraction of clean tiles with a false detection. (b) Layout network: intersection over union (IoU) per class.}
\label{fig:training}
\end{figure*}

\section{Experimental Results}\label{sec:results}
All experiments use superconducting quantum circuits as the test vehicle, scanned in brightfield on an Olympus DSX1000. The set comprises two 100-mm wafers of one design (NFT1 and 009-006~\#2, each with three reticle designs placed 12 times, at 1.63\umpx), chip AHQ1 A2 (1.0\umpx), and chips A2, B2 and C1 of the AFT5 series. Chip C1 is excluded from all training. The risk results use the superconducting instance of Table~\ref{tab:risk}; every other stage runs unchanged from its process-agnostic form.

\subsection{Planted-Defect Benchmark}
Without ground-truth defect maps for real scans, detection is measured by planting. From the held-out chip C1, 150 windows of $768 \times 768$~px are sampled. Each window is processed twice by the complete pipeline, verification chain included: once as scanned, and once with two synthetic specks planted, 300 in total. A planted speck counts as found when a reported finding overlaps it, so a miss at any stage is a miss. Detections on the unplanted copies give the false-alarm rate, an upper bound because the clean windows contain real dust.

Fig.~\ref{fig:recall} breaks recall down by speck area and contrast for three generations of the speck network. Specks above 60~px were already found about 96\% of the time by speck\_v2; the later gains are concentrated in small and faint specks. Recall on specks under 12~px rose from 64.4\% to 86.7\%, and on specks with contrast below 0.06 from 65.7\% to 80.0\%. Overall recall is 89.0\%, 92.3\% and 93.3\% for speck\_v2, speck\_v3 and speck\_v4.

\begin{figure*}[!t]
\centering
\includegraphics[width=\textwidth]{fig3_recall.png}
\caption{Recall of the full pipeline on 300 specks planted into held-out chip C1, by speck network. (a) By speck area. (b) By contrast against the local background. Group sizes are given under each axis label.}
\label{fig:recall}
\end{figure*}

Fig.~\ref{fig:tradeoff} places each configuration on a recall--false-alarm plane. Each gain in recall was bought with more detections on clean windows: from 0.23 per window for speck\_v2 to 0.63 for speck\_v4. The final point is not a new network but the context-conditioned verification of Section~\ref{sec:verification}. Recall is unchanged at 93.3\% while clean-window detections rise to 1.21, because isolated specks on open ground are no longer rejected. Inspection of the additional clean-window detections from speck\_v4 showed most to be real dust in the scan rather than false alarms.

\begin{figure}[!t]
\centering
\includegraphics[width=\columnwidth]{fig4_tradeoff.png}
\caption{Recall versus detections per clean $768 \times 768$~px window on chip C1, for successive pipeline configurations.}
\label{fig:tradeoff}
\end{figure}

\subsection{Wafer-Scale Results}
Table~\ref{tab:findings} lists the findings on three production scans. Context-conditioned verification had the largest effect on the wafers. Before it, the lettering and device checks rejected 94{,}600 of 123{,}711 candidates on wafer 009-006~\#2; random samples of those rejections were almost all real particles, fibres and faint dust on open ground. With the change, NFT1 findings rose from 27{,}133 to 119{,}122. Most of the increase was low-risk dirt, but critical findings also rose from 377 to 449 and high from 571 to 995.

\begin{table*}[!t]
\caption{Findings per Scan With speck\_v4 and Context-Conditioned Verification}
\label{tab:findings}
\centering
\begin{tabular}{@{}lrrrrrrr@{}}
\toprule
Scan & Size (px) & Pixel (\textmu m) & Critical & High & Medium & Low & Total \\
\midrule
NFT1 wafer & $118{,}092 \times 118{,}088$ & 1.63 & 449 & 995 & 620 & 117{,}058 & 119{,}122 \\
009-006 \#2 wafer & $118{,}092 \times 118{,}088$ & 1.63 & 354 & 326 & 590 & 116{,}779 & 118{,}049 \\
AHQ1 chip A2 & $44{,}627 \times 50{,}500$ & 1.00 & 37 & 64 & 176 & 4{,}917 & 5{,}194 \\
\bottomrule
\end{tabular}
\end{table*}

At wafer scale, individual outlines are unreadable, so the review interface aggregates findings into a density map over 0.5~mm cells (Fig.~\ref{fig:heatmaps}). On NFT1 the map separates three patterns: findings concentrated along the wafer edge and reticle borders; a broad diagonal band across the lower half; and a patch on the left. The band disappears when only medium-risk and higher findings are counted, which identifies it as low-risk surface contamination; the serious findings are spread nearly uniformly over the reticles and crowd at the wafer edge.

\begin{figure*}[!t]
\centering
\includegraphics[width=\textwidth]{fig5_heatmaps.png}
\caption{Density of findings on the NFT1 wafer in 0.5~mm cells, drawn on the chip silhouette with the image hidden. (a) All 119{,}122 findings. (b) The 2{,}064 findings rated medium risk or higher.}
\label{fig:heatmaps}
\end{figure*}

\subsection{CPW Metrology}
On NFT1, 16{,}222 sites were measured; wafer 009-006~\#2 was registered by initializing from the NFT1 transform, in 3.5~min, and gave 17{,}831 sites. Fig.~\ref{fig:metrology} shows the mean deviation from nominal for every structure width measured at least 30 times. The two wafers agree to within 0.1\um{} on the well-sampled widths: nominal 10\um{} gaps measure $8.76 \pm 0.34$\um{} and $8.75 \pm 0.33$\um{} (mean $\pm$ standard deviation, $n = 7{,}228$ and $7{,}478$). Lines measure wider and gaps narrower by comparable amounts ($+0.75$\um{} for 100\um{} lines, $-1.43$\um{} for 100\um{} gaps), consistent with each metal edge imaging 0.4--0.6\um{} outward of its designed position. The agreement establishes that the offset is reproducible; whether it originates in lithography or in where the dark edge line forms in brightfield requires a cross-check by electron microscopy. Widths of 4\um{} and below (under 3~px) are not resolved and are excluded.

\begin{figure*}[!t]
\centering
\includegraphics[width=\textwidth]{fig6_metrology.png}
\caption{Mean measured minus nominal width for CPW lines (metal) and gaps on two wafers of the same design, for widths measured at least 30 times. Error bars are one standard deviation.}
\label{fig:metrology}
\end{figure*}

\subsection{Device-to-Device Comparison}
A component check was also evaluated, which compares every instance of a repeated device, such as a junction or a resonator coupler, with a reference built from other instances. With the reference taken as the median of instances in the same row, review judged essentially all flags on the 009-006 wafers and AHQ1 A2 to be wrong; flagged structures matched real neighbours at normalized cross-correlation 0.87--0.98. Rows mix device types that share a footprint, so the median was the wrong reference. To quantify this, 36 crude defects (particles, bridges and missing metal) were painted onto 12 device types of AHQ1 A2. Fig.~\ref{fig:components} compares the row-median reference with one built from each device's 12 most similar instances, across a range of thresholds. The look-alike reference finds 27 of 36 at 54 flags among 67{,}448 clean devices, against 6 of 36 at 719 flags for the row median. Its remaining flags were mostly pad corners, junction-lead edges and stitching double images, so the check is disabled by default. Dirt on devices remains covered by the speck detectors.

\begin{figure}[!t]
\centering
\includegraphics[width=\columnwidth]{fig7_components.png}
\caption{Planted device defects found (of 36) versus flags on the clean scan (67{,}448 devices) for the row-median reference and for the look-alike reference at thresholds $z = 4$ to 20.}
\label{fig:components}
\end{figure}

\subsection{Runtime}
The NFT1 wafer, 9{,}264 network tiles, was processed in 37~min on an Apple M5 laptop: 17~min of network inference and 20~min of verification and risk scoring. Parallelizing the verification chain across CPU cores accelerated it 9.3-fold with identical output. Training speck\_v4 took 54~min on the workstation GPU.

\section{Discussion and Limitations}\label{sec:discussion}
\emph{Generality.} The pipeline makes three assumptions about a substrate: that it is patterned, that structures repeat somewhere on the scan (for the reference residual and twin verification), and that defects contrast with their surroundings in the chosen illumination. It makes none about materials, line widths or device types. Moving to a new process means re-running the label-free training on that process's scans and supplying its risk table; no annotation is required. The evidence for this generality is limited to one microscope and one device family; the training set also included printed-circuit scans, but no detection results are reported for them. Validation on CMOS, MEMS or photonic wafers, and on darkfield or reflected-light differential-contrast imaging, is the most important next experiment.

\emph{Benchmark validity.} Planted specks are drawn from the same family of shapes used in training, so the benchmark measures recall on that family. Real defects outside it (hairs, stains, lifted films) are covered by the anomaly network, but its recall on them is not measured here. The clean-window rate counts real dust as false alarms and therefore overstates them. Both limitations would be removed by a set of reviewed real findings; the interface already records reviewer verdicts, 266 to date across three scans, as a basis for one.

\emph{Context conditioning.} Rejection rules that depend on a learned label inherit that label's errors. The layout network labels isolated particles as lettering or device structure, and a rule that trusted the label alone discarded three quarters of all candidates on one wafer. Conditioning the rule on local context removed that failure without changing benchmark recall. This suggests that recall on planted defects alone does not catch every failure of the verification chain.

\emph{Resolution.} At 1.6\umpx, structures below about 5\um{} (3~px) cannot be measured, and defects of one or two pixels are near the detection floor; recall on specks under 12~px is 86.7\%. About 10\% of wafer registration control points still deviate by more than 20~px, which affects metrology locally but not detection, which is design-free.

\emph{Scope.} Device-to-device comparison remains disabled pending labelled device damage. Comparison against the GDS for extra or missing metal is a prototype that requires per-layer rendering and an alignment robust to repeated structures. Darkfield scans, which image particles at higher contrast, are not yet used in training.

\section{Conclusion}\label{sec:conclusion}
Chip Scanner turns a gigapixel brightfield scan of a patterned wafer into a ranked list of risk-scored findings in 37~min on a laptop, without manual labels and without assuming a device technology outside its risk table. Its detectors are trained on defects injected into the scans themselves, so it adapts to a new process from that process's images. On superconducting quantum circuits, used here as a demanding test vehicle, it recovers 93.3\% of planted specks on a held-out chip through the complete pipeline. Its context-conditioned verification keeps isolated particles that layout labels alone would discard. When the design is available, the same scan yields sub-pixel line and gap metrology at more than 16{,}000 sites, reproducible to 0.1\um{} across two wafers. The next steps are validation on other technologies and illumination modes, a reviewed set of real findings to replace planted-defect evaluation, and a validated device-comparison check.

\bibliographystyle{IEEEtran}
\bibliography{refs}

\end{document}
